% % Anselmo/Nick
% % fundamental-concepts-phenotypic-plasticity.tex
% %%%%%%%%%%%%%%%%%%%%%%%%%%%%%%%%%%%%%%%%%%%%%%%%%%%%%%%%%%%%%%%%%%%%%%%%%%%%%%%%%%
% % revision status as of March 3rd 2026
% % - [ ] accept or reject all track changes
% % - [ ] incorporate review comments
% %    - don't worry about clearing all comments
% % - [ ] incorporate reference suggestions
% % - [ ] cross-cutting revision objectives (see email)
% % - [ ] other planned changes or refinements (if any)
% % for revision status, mark - for partial or x for complete
% % any free-form comments on planned/completed revisions are helpful too!
% % - Will revise rough draft, shorten it and add citations.
% % - The discussion on learning will change substantially.
% % --------------------------------------------------------------------------------
% % draft status: <rough draftt> as of <03-Mar-2026>
% % pick one of: outline, partial draft, rough draft, complete draft
% % where rough draft indicates that major points are hit,
% % but still working on changes or improvements in refining the text
% %%%%%%%%%%%%%%%%%%%%%%%%%%%%%%%%%%%%%%%%%%%%%%%%%%%%%%%%%%%%%%%%%%%%%%%%%%%%%%%%%%
% \begin{bibunit}

% % please ensure refs have a DOI where possible!!!
% \begin{filecontents*}[overwrite]{fundamental-concepts-phenotypic-plasticity.bib}
% @book{Darwin2009,
%   title = {The Origin of Species: By Means of Natural Selection,  or the Preservation of Favoured Races in the Struggle for Life},
%   ISBN = {9780511694295},
%   url = {http://dx.doi.org/10.1017/CBO9780511694295},
%   DOI = {10.1017/cbo9780511694295},
%   publisher = {Cambridge University Press},
%   author = {Darwin,  Charles},
%   year = {2009},
%   month = jul
% }
% \end{filecontents*}

\subsection{Phenotypic Plasticity}
\label{sec:phenotypic-plasticity}

% \textbf{Suggested Citations}: \begin{itemize}
% \item find full reference info in \texttt{filecontents} above, or in rendered PDF
% \item empty this list by moving citations into ``Included Citations'' or ``Excluded Citations'' below
% \end{itemize}

% \textbf{Included Citations}: \begin{itemize}
% \item move citations incorporated into section text here!
% \end{itemize}

% \textbf{Excluded Citations}: \begin{itemize}
% \item move citations not incorporated into section text here!
% \end{itemize}

Phenotypic plasticity is the ability of organisms to modify their phenotypes in response to environmental conditions.
In simpler terms, it means the ability that living beings have to adjust to their environment during their lifetime,  altering not only form or behavior but also development, physiology, and other traits to cope with changing circumstances.
It is a universal feature of life, from bacteria to animals and plants, and represents one of nature's most powerful adaptive mechanisms.
It is also the characteristic most associated with intelligence.
While every regulated process involves a degree of intelligence, phenotypic plasticity encompasses the most conspicuous expressions of it, such as learning.
For evolutionary computation (EC), which seeks to reproduce the adaptive power of evolution in silico, phenotypic plasticity offers a set of principles that can expand flexibility, generality, robustness, and long-term evolvability of artificial systems.

In biology, plasticity appears in two main forms: developmental and behavioral.
Developmental plasticity involves structural or physiological adjustments that occur during growth, such as grasshoppers changing color or Daphnia developing protective spines when predators are present.
Behavioral plasticity, by contrast, is rapid and reversible, allowing organisms to modify actions within their lifetime.
Examples include bacteria adjusting metabolic pathways to different nutrients or animals learning to navigate, hunt, or communicate.
These forms of plasticity depend not only on immediate sensory inputs but also on stored information, biological memory built from past experiences.
Learning, a specific case of behavioral plasticity, enables individuals to refine responses over time and represents a key evolutionary innovation in intelligent organisms.

The mechanisms underlying plasticity span molecular, cellular, and social scales.
Epigenetic mechanisms such as DNA methylation and histone modification allow gene expression to be tuned by environmental inputs and to persist through cell divisions.
These chemical marks provide a form of semi-heritable memory that helps organisms ``remember'' previous conditions.
Neurological plasticity operates at the level of neurons and synapses, enabling changes in connectivity and strength through processes like long-term potentiation and neuromodulation, biological foundations of learning and memory.
At a larger scale, populations also store and transmit adaptive information through horizontal gene transfer or cultural inheritance, extending memory beyond the individual.
Memory is central to all these forms of plasticity: it allows organisms to adjust not only reactively but anticipatorily, based on prior states.
This concept maps naturally onto computational analogs.
Epigenetic memory parallels adaptive parameter tuning or structural modification in EC systems.
Neurological plasticity corresponds to weight adaptation or dynamic network topology.
Population-level memory resembles knowledge sharing among agents or runs in distributed evolutionary algorithms.

Despite its ubiquity in nature, phenotypic plasticity has only limited representation in EC.
Most systems maintain a fixed genotype-phenotype mapping, evolving static solutions that lack the ability to change during their lifetime.
While some studies explore related ideas, the Baldwin effect, Lamarckian learning, or adaptive operators, these remain peripheral rather than central design principles.
As a result, EC often captures evolution's selective power but not its developmental or learning dimensions.

Integrating richer forms of plasticity into EC could bridge evolution and learning in artificial systems.
In genetic programming, for instance, evolving programs that can adapt or learn during execution would move closer to biological intelligence.
Even minimal internal memory, such as a few mutable registers, self-modifying code, or feedback loops, can enable adaptive behaviors to emerge.
Incorporating biologically inspired mechanisms that mimic epigenetic regulation or learning processes promises to enhance robustness, generalization, and long-term open-ended evolution in EC systems.
Phenotypic plasticity, long central to biology, may hold the key to unlocking new levels of adaptability in artificial evolution.


% \putbib[fundamental-concepts-phenotypic-plasticity]

% \end{bibunit}
